\begin{tuluyen}
Trong nỗ lực giảm phụ thuộc vào nhiên liệu hóa thạch và bảo vệ khí quyển, Việt Nam đã thay thế xăng truyền thống bằng xăng sinh học E10. Xăng E10 là loại xăng sinh học chứa 10\% ethanol về thể tích ($D = 0,789$ g/mL), còn lại là xăng truyền thống (giả thiết chỉ chứa hai alkane là \ce{C8H18} và \ce{C9H20} có tỉ lệ mol tương ứng là 3:4 và $D = 0,7$ g/mL).

Một người đi xe máy di chuyển từ phường Xuân Hương -- Đà Lạt đến phường Phan Thiết với quãng đường 180 km. Trong bình xăng của xe có 3,0 L xăng E10. Nhiệt lượng tỏa ra khi đốt cháy hoàn toàn 1 mol các chất như bảng sau:
\begin{center}
\begin{tabular}{c|ccc}
\toprule
Chất & \ce{C2H5OH} & \ce{C8H18} & \ce{C9H20} \\
\midrule
Nhiệt lượng tỏa ra (kJ/mol) & 1 375 & 5 095 & 6 115 \\
\bottomrule
\end{tabular}
\end{center}
Biết rằng để di chuyển được 1 km, động cơ xe máy cần một nhiệt lượng để chuyển thành công cơ học có độ lớn là 210 kJ và hiệu suất sử dụng nhiên liệu của động cơ là 32\%.

Hãy tính và cho biết để đi hết quãng đường trên thì người đó có cần đổ thêm xăng không?
\source{Trích nguyên văn Câu 2.2, đề tuyển sinh chuyên Hóa Lâm Đồng 2026.}
\dapso{Với 3,0 L E10 xe đi được khoảng \SI{144.8}{\kilo\metre}; cần đổ thêm khoảng \SI{0.73}{\liter} E10 để đi đủ \SI{180}{\kilo\metre}.}
\begin{loigiaiBT}
Thể tích ethanol:
\[
V_{\ce{C2H5OH}}=3.0\times10\%=\SI{0.30}{\liter}=\SI{300}{\milli\liter}.
\]
\[
m_{\ce{C2H5OH}}=300\times0.789=\SI{236.7}{\gram},\qquad
n_{\ce{C2H5OH}}=\frac{236.7}{46}\approx\SI{5.146}{\mole}.
\]
Phần xăng truyền thống có thể tích \SI{2.70}{\liter}, khối lượng:
\[
m=2700\times0.7=\SI{1890}{\gram}.
\]
Đặt $n_{\ce{C8H18}}=3x$, $n_{\ce{C9H20}}=4x$:
\[
3x\times114+4x\times128=1890\Rightarrow x\approx2.213.
\]
Suy ra:
\[
n_{\ce{C8H18}}\approx\SI{6.639}{\mole},\qquad
n_{\ce{C9H20}}\approx\SI{8.852}{\mole}.
\]
Tổng nhiệt lượng tỏa ra:
\[
Q\approx5.146\times1375+6.639\times5095+8.852\times6115
\approx\SI{95036}{\kilo\joule}.
\]
Công cơ học có ích:
\[
Q_{\text{ích}}=95036\times0.32\approx\SI{30411}{\kilo\joule}.
\]
Quãng đường đi được:
\[
s=\frac{30411}{210}\approx\SI{144.8}{\kilo\metre}.
\]
Do $144.8<180$, người đó phải đổ thêm xăng. Vì cùng loại E10, thể tích cần thêm tỉ lệ với quãng đường:
\[
V_{\text{thêm}}=3.0\left(\frac{180}{144.8}-1\right)\approx\SI{0.73}{\liter}.
\]
\textbf{Vậy,} người đó cần đổ thêm xăng; lượng E10 tối thiểu cần thêm khoảng \SI{0.73}{\liter}.
\end{loigiaiBT}

\end{tuluyen}
